\begin{abstract}
Machine-learning research is increasingly central to scientific discovery, yet many studies remain difficult to evaluate, reproduce, or translate into defensible publication claims. Existing reporting standards define important disclosure expectations, but they do not by themselves provide an operational workflow for producing those disclosures from the beginning of the research process. This article proposes a query-to-publication scaffold for reproducible machine-learning science. The scaffold organizes research into staged, version-controlled artifacts: research query, study protocol, data audit, feature plan, modeling plan, validation plan, result bundle, ethics review, reproducibility package, manuscript drafts, and submission materials. Its central contribution is an evidence-to-claim control model that links manuscript claims to supporting artifacts and bounds those claims by the strength of validation. The framework positions reproducibility, reviewability, and claim discipline as properties of the full scientific workflow rather than as final-stage manuscript additions. By operationalizing the path between reporting standards and project-specific code, the scaffold offers a practical workflow layer for researchers, reviewers, editors, and instructors seeking more reliable machine-learning manuscripts.
\end{abstract}
